\section{Introduction}

Cooperative UAV decision making is moving from platform-centric pursuit toward networked sensing, communication, and engagement. In such missions, success is not determined only by whether a single aircraft approaches a target. A team must detect the target, transfer or relay target information, keep usable communication paths, position an appropriate shooter, and form an engagement window before the target escapes or the team loses situational awareness. This sequence is treated in this paper as a cooperative kill chain.

Heterogeneous UAV teams make this problem more difficult. A scout UAV may have a wider radar field of view but weaker attack capability. A relay UAV may have a larger communication range but weaker maneuverability. An attacker may have better engagement geometry but may depend on target information discovered by another platform. Therefore, the relation between two UAVs is not simply a generic graph edge. A scout-to-attacker edge can represent target-information support, a relay-to-attacker edge can represent communication reconstruction, and an attacker-to-relay edge can represent the need to preserve a useful information path.

This paper studies an implementable intermediate step toward high-fidelity air-combat decision making. Instead of directly training all baselines in a full 6DOF JSBSim environment, we build a lightweight constrained 3DOF tactical interception environment. It preserves the key decision constraints needed for the present claim: three-dimensional position, speed, heading, flight-path angle, altitude limits, turn and climb constraints, radar detection, limited communication, message age, and attack-window metrics.

The proposed method is EA-RG-MAPPO-S, a task-graph-driven multi-relation MAPPO method for post-failure coordination recovery under limited communication. It represents the UAV team and target as a multi-relation task graph, separates perception, communication, and task-support relations, and processes them with state-dependent, edge-feature-modulated attention. A structured role-pair initialization prior (Gate Prior) stabilizes role-structured optimization, and a static role-pair modulation is retained as an auxiliary architectural component. Under strict intermittent sensing and relay-node communication failure, a pre-specified censor-aware survival analysis (RQ2) shows that the most reproducible benefit is \emph{earlier} post-failure recovery relative to MAPPO: the primary restricted mean survival time is lower than MAPPO and the wider single-graph baseline, while remaining competitive with HAPPO over the full horizon. The method therefore shifts the temporal concentration of post-failure recovery rather than claiming a uniformly faster recovery-time distribution.

The contributions are:
\begin{enumerate}
    \item We formulate a strict-sensing 3DOF heterogeneous UAV kill-chain recovery task with relay-node communication failure and post-failure recovery latency as a first-class metric.
    \item We propose a task-graph-driven multi-relation graph policy (perception / communication / task-support relations with state-dependent edge-feature-modulated attention) for decentralized actors, in which task-support acts as a task-dependent relational mask over delivered communication.
    \item We provide a locked three-seed evidence chain (10{,}800-episode held-out, 10{,}500-episode robustness, efficiency profiling) with a censor-aware survival analysis showing near-saturated recovery reliability and a reproducible \emph{early-recovery} advantage relative to MAPPO (competitive with HAPPO over the full horizon), plus mechanism evidence that Gate Prior improves cross-seed optimization stability.
\end{enumerate}
